\documentclass[12pt]{article}
\usepackage[margin=0.58in,top=0.62in,bottom=0.48in]{geometry}
\usepackage{lmodern}
\usepackage{microtype}
\usepackage{fancyhdr}
\usepackage{titlesec}
\usepackage{enumitem}
\usepackage{lastpage}
\usepackage[hidelinks]{hyperref}

\setlength{\parindent}{0pt}
\setlength{\parskip}{0.18em}
\setlength{\headheight}{26pt}
\raggedbottom
\titleformat{\section}{\normalsize\bfseries\scshape}{}{0pt}{}
\titleformat{\subsection}{\normalsize\bfseries}{}{0pt}{}
\titlespacing*{\section}{0pt}{0.35em}{0.12em}
\titlespacing*{\subsection}{0pt}{0.25em}{0.08em}
\pagestyle{fancy}
\fancyhf{}
\fancyhead[C]{\footnotesize Lit Example Reader \quad Assignment ID: U1L15LE \quad Page \thepage\ of \pageref{LastPage}}
\renewcommand{\headrulewidth}{0.5pt}

\newenvironment{playpassage}{\begin{quote}\footnotesize\setlength{\parskip}{0.08em}\setlength{\parindent}{0pt}}{\end{quote}}
\newcommand{\speaker}[1]{\textbf{#1.}}

\begin{document}

\begin{center}
{\LARGE\bfseries Week 15 Lit Example Reader}\par\vspace{0.02em}
{\large\scshape Shakespeare: Words That Carry Disputes}\par\vspace{0.06em}
\rule{0.78\textwidth}{0.5pt}
\end{center}

\textbf{Purpose:} Shakespeare passages for Week 15: decide whether a conflict is verbal, real, or mixed, and name what remains disputed after key words are clarified.

\textbf{What to watch for:} Speakers often use powerful words such as \textit{man}, \textit{honor}, \textit{ambition}, or \textit{liberty}. Ask what each speaker means and what would still need proof after the word is clarified.

\section*{Before the Passages: The Week 15 Logic Tool}

\begin{itemize}[leftmargin=1.3em,itemsep=0.05em,topsep=0.08em,parsep=0.03em]
\item \textbf{Verbal dispute:} mostly a disagreement over word meaning.
\item \textbf{Real dispute:} remains after words are clarified.
\item \textbf{Mixed dispute:} word-confusion plus a real disagreement underneath.
\item \textbf{Dispute map:} contested term; Side A meaning; Side B meaning; shared claim; remaining disagreement; classification.
\end{itemize}

\section*{Passage 1: Macbeth, Act 1, Scene 7 --- What Does It Mean to Be a Man?}

\textbf{Context:} Macbeth hesitates before murdering Duncan. Lady Macbeth attacks his manhood. Macbeth answers with a different account of what a man should dare.

\begin{playpassage}
\speaker{Lady Macbeth} When you durst do it, then you were a man;\\
And, to be more than what you were, you would\\
Be so much more the man.\\
...\\
\speaker{Macbeth} Prithee, peace:\\
I dare do all that may become a man;\\
Who dares do more is none.
\end{playpassage}

\subsection*{Reader Notes}

This is your assisted example. The contested term is \textit{man}. Lady Macbeth treats manhood as boldness that will carry out the murder. Macbeth's reply points toward a different meaning: a true man dares what is fitting, but not what violates the moral boundary of being human. Clarifying the word helps, but the dispute is not merely verbal. They still disagree about whether murder is courage, excess, loyalty, ambition, or moral failure.

\textbf{Reading task:} Circle each use of \textit{man}. In the margin, write Lady Macbeth's meaning and Macbeth's meaning in five words or fewer each.

\newpage

\section*{Passage 2: Julius Caesar, Act 3, Scene 2 --- Honor, Ambition, and the Real Point}

\textbf{Context:} Brutus tells Rome that Caesar died because he was ambitious. Antony repeatedly calls Brutus and the conspirators honorable while presenting evidence that presses against Brutus's claim.

\begin{playpassage}
\speaker{Brutus} As Caesar loved me, I weep for him; as he was fortunate, I rejoice at it; as he was valiant, I honour him; but, as he was ambitious, I slew him.\\
...\\
\speaker{Antony} The noble Brutus\\
Hath told you Caesar was ambitious:\\
If it were so, it was a grievous fault,\\
And grievously hath Caesar answer'd it.\\
...\\
He hath brought many captives home to Rome\\
Whose ransoms did the general coffers fill:\\
Did this in Caesar seem ambitious?\\
When that the poor have cried, Caesar hath wept:\\
Ambition should be made of sterner stuff.
\end{playpassage}

\subsection*{Reader Notes}

This passage shifts more work to you. The words \textit{honorable} and \textit{ambitious} carry the public dispute. Do not stop at noticing that Antony is ironic. Ask what Brutus needs \textit{ambitious} to mean, what evidence Antony offers against that meaning, and whether clarifying the word would end the political conflict.

\textbf{Reading task:} Underline Brutus's claim about Caesar. Box Antony's examples. Put a star beside one repeated word whose meaning is being pressured.

\section*{How to Use These Passages on the Worksheet}

Do not summarize the speeches. Use Week 15's dispute map: contested term, each side's meaning, shared claim, remaining disagreement, and verbal/real/mixed classification.

\end{document}
